% !TeX root = ../深度学习与大语言模型技术原理.tex
% 由主文档通过 \input 载入；请编译主文档。
\setcounter{chapter}{5} % 保留原书章号。
\chapter{专家混合、稀疏计算与混合架构}
\label{ch:moe}

\section{条件计算与路由}

稠密 FFN 对每个词元使用同一组参数。\emph{专家混合（mixture of experts, MoE）}保留多组专家网络 $E_1,\ldots,E_K$，让路由器（router）根据输入选择其中少数几个。设 $p(x)=\operatorname{softmax}(W_rx)$，$S_k(x)$ 为概率最大的 $k$ 个专家索引，可定义
\begin{equation}
y=\sum_{e\in S_k(x)}g_e(x)E_e(x),\qquad
g_e(x)=\frac{p_e(x)}{\sum_{j\in S_k(x)}p_j(x)}.
\label{eq:moe}
\end{equation}
这是重新归一化选中权重的一种约定；也有实现直接保留原始门控概率。Top-$k$ 选择是离散操作，训练通常沿已选分支的连续权重传播梯度，不能把整个选择过程当成处处可微函数。

MoE 增加总参数容量，而每词元只执行少量专家。若每个专家有 $P_e$ 个参数，共享部分有 $P_s$ 个参数，则总参数近似为 $P_s+KP_e$，每词元激活参数近似为 $P_s+kP_e$。后者有助于估算算术成本，前者影响权重存储与加载；通信、路由和填充又会增加额外成本。因此，相同激活参数量不保证相同延迟。

\section{容量、负载均衡与训练失败}

一个含 $N$ 个词元的路由批次会产生约 $kN$ 次专家分配。给每个专家容量 $C=\lceil c kN/K\rceil$，其中 $c$ 是容量系数，可以限制单专家缓存峰值。超容量词元需要丢弃、重路由或采用不丢词元的动态分配；这些是不同算法选择，必须记录，尤其不能悄悄改变有效训练样本。

为说明负载均衡，考虑每词元只选一个专家。令 $f_e$ 为被分配给专家 $e$ 的词元比例，$\bar p_e$ 为该专家门控概率的批平均值，可加入
\[
\mathcal L_{\mathrm{bal}}=\alpha K\sum_{e=1}^Kf_e\bar p_e.
\]
通常不通过离散计数 $f_e$ 反向传播，而通过 $\bar p_e$ 推动概率分配。它在实践中用于抑制路由拥挤，不能单凭这个乘积形式宣称每个批次必然达到均匀负载。Switch Transformer 给出了此类稀疏路由与辅助损失的具体实例\cite{switch}。

若少数专家长期获得大多数词元，其他专家缺少学习机会，可能出现专家坍缩（expert collapse）。反之，过强均衡又可能压制自然的任务分工。应同时监测每专家词元数、路由熵、溢出率、专家梯度和主任务损失；只看总体损失不能定位路由异常。

\section{专家并行与通信路径}

专家并行（expert parallelism, EP）把专家放到不同设备，路由后把词元表示发送给目标设备，计算完成后再返回并加权合并。该过程常包含 All-to-All 通信，传输对象主要是激活及其反向梯度，而不只是路由编号。各设备专家负载不同会产生等待，跨节点带宽和通信拓扑可能主导速度。

共享专家可处理多数词元共用的知识，路由专家承担更专门的变换；多个专家也可被同一设备持有。推理时请求数量少会使专家矩阵乘法批量过小，降低硬件利用率。因此 MoE 训练与服务必须共同考虑，不能只由理论 FLOPs 判断部署成本。它与数据、张量和流水线并行的组合见第~\ref{ch:distributed} 章。

\section{稀疏注意力、状态空间与混合模型}

MoE 稀疏的是参数激活，稀疏注意力稀疏的是位置之间的连接。若每个词元只访问窗口内 $w$ 个位置，注意力计算可从 $O(n^2d)$ 降为 $O(nwd)$，但远距离信息需要通过多层传播或额外全局位置传递。除非被删去的注意力权重原本为零，否则它通常改变模型函数。

\emph{状态空间模型（state space model, SSM）}用递推状态汇总历史。一个离散线性形式为 $s_t=As_{t-1}+Bx_t$、$y_t=Cs_t+Dx_t$。若状态维度固定，逐词元更新无需缓存全部历史键值，但状态必须压缩过去信息。选择性状态空间模型使部分转移或读写参数依赖输入，以增强内容选择能力；Mamba 是这一思路的具体架构\cite{mamba}，不能把上面的固定线性递推当成其完整实现。

混合架构（hybrid architecture）交替使用注意力、递推模块或专家层，在精确检索历史内容与有限状态记忆之间折中。比较时应固定训练数据和预算，并在长距离检索、顺序建模、吞吐量及状态存储上分别评估。线性序列复杂度是系统优势的一部分，却不证明对所有语言任务都有更高准确率。

